\documentclass[11pt]{article}
\usepackage[margin=1in]{geometry}
\usepackage{amsmath,amssymb,amsthm}
\usepackage{microtype}
\usepackage{parskip}

\title{Can the $n=2$ response function correct interference?}
\author{}
\date{}

\newcommand{\tr}{\operatorname{tr}}
\newcommand{\RR}{\mathbb{R}}
\newcommand{\EE}{\mathbb{E}}

\begin{document}
\maketitle

\section{Setup and question}

A model has response $R(v) = F(S(v))$ with
\[
  S(v) = \sum_{i=1}^M w_i (d_i^\top v)^2 = v^\top A v,
  \qquad A = \sum_i w_i d_i d_i^\top,
\]
unit feature directions $\{d_i\}$ in $\RR^D$ with $M > D$, weights
$w_i \ge 0$, and plateau $\{v : S(v) < c\}$. The plateau-breaking
magnitude along unit $d$ is $m_d = \sqrt{c / (d^\top A d)}$. Write
$q_i := d_i^\top A d_i$.

With $k$ active features, an inactive feature $i$ picks up spurious
read-off of typical size $\sigma_a\sqrt{k/D}$. We demand:

\begin{itemize}
  \item \textbf{Per-feature correction:}
    $\sigma_a\sqrt{k/D} \le m_{d_i} \le \sigma_a$ for every $i$ ---
    i.e.\ $c/\sigma_a^2 \le q_i \le Dc/(k\sigma_a^2)$.
  \item \textbf{LRH sensitivity:} $m_{d_i} < m_{d_r}$ for every
    feature $i$ (random unit $d_r$) --- i.e.\ $q_i > \tr(A)/D$.
\end{itemize}

\textbf{Question.} Do weights $\{w_i\}$ and $c > 0$ exist making both
conditions hold simultaneously?

\section{Answer (i): under tight framing, no}

If the feature frame is tight ($\Phi\Phi^\top = (M/D)\,I$), the answer
is \textbf{no, for any weight choice}.

Summing LRH over features:
\[
  \sum_i q_i = \tr(A\,\Phi\Phi^\top) = \frac{M}{D}\,\tr A.
\]
So the average of $q_i$ over features equals $\tr(A)/D$ exactly --- at
least one feature must violate strict LRH, regardless of
$\{w_i\}$. No reweighting rescues a tight frame.

\paragraph{Why tight framing is the usual assumption.}
Tight frames minimise worst-case coherence $\max_{i\ne j}|G_{ij}|$
(Welch bound), equalise the spurious read-offs $\hat a_i$ across
features, and make $\Phi\Phi^\top$ isotropic so that random directions
in activation space are all treated the same. These are exactly the
conditions under which the typical-interference estimate
$\sigma_a\sqrt{k/D}$ is valid. Tightness is therefore the natural
idealisation whenever one wants interference to be uniformly small
across features --- and the $n=2$ response function cannot
simultaneously satisfy interference correction and LRH sensitivity in
this idealisation.

\section{Answer (ii): lifting tight framing, yes}

If the frame is not tight, the answer is \textbf{yes}. The simplest
construction: take directions with $G_{ij}^2 \approx 1/D$ for $i \ne
j$ (i.e.\ the regime in which the $\sigma_a\sqrt{k/D}$ estimate is
accurate, but without enforcing $\Phi\Phi^\top = (M/D)I$). With
uniform weights $w_i = 1$,
\[
  q_i = 1 + \sum_{j\ne i} G_{ij}^2 \;\approx\; \frac{M}{D} + \frac{D-1}{D}.
\]
Then $q_i - \tr(A)/D \approx (D-1)/D$: a margin of order $1$ per
feature, so LRH holds for every $i$. Choose $c = \sigma_a^2 \min_i
q_i$: the feature with smallest $q_i$ hits $m = \sigma_a$; all others
satisfy $q_i \in [q_{\min},\, q_{\max}]$ with $q_{\max}/q_{\min} = 1 +
O(\sqrt{M}/D)$, easily inside the factor-$D/k$ bracket for $k \ll D$.

\paragraph{What this costs.}
Dropping tight framing buys feasibility but introduces a parasitic
privileged direction. The top eigenvector $u$ of $\Phi\Phi^\top$ has
$q_u = \lambda_{\max}(A) > q_i$ for every feature, so $u$ --- which is
\emph{not} a feature direction --- is steeper than any feature. LRH
then holds only in the typical sense (features beat typical random
directions), not universally (features do not beat the worst-case
random direction). Non-uniform interference also appears: feature
pairs with $G_{ij}^2$ above average exceed the Welch bound, and the
$\sigma_a\sqrt{k/D}$ estimate becomes an approximation rather than a
uniform bound.

\section{Summary}

\begin{itemize}
  \item Tight frame $\Rightarrow$ LRH and per-feature correction
    cannot coexist, for any weights.
  \item Non-tight frame with $G_{ij}^2 \approx 1/D$ uniformly
    $\Rightarrow$ both conditions hold with uniform weights.
  \item The cost of dropping tightness is a non-feature direction
    (the top eigenvector of $\Phi\Phi^\top$) that is strictly steeper
    than any feature.
\end{itemize}

\end{document}
